% 海森伯群（综述）
% license CCBYSA3
% type Sum

本文根据 CC-BY-SA 协议转载翻译自维基百科\href{https://en.wikipedia.org/wiki/Heisenberg_group}{相关文章}

在数学中，海森堡群$H$（以维尔纳·海森堡命名）是由如下形式的$3 \times 3$ 上三角矩阵组成的群：
$$
\begin{pmatrix}
1 & a & c \\
0 & 1 & b \\
0 & 0 & 1
\end{pmatrix}~
$$
其运算为矩阵乘法。元素$a, b, c$可以取自任意含幺交换环，常见的选择是实数环（得到所谓的“连续海森堡群”）或整数环（得到“离散海森堡群”）。

连续海森堡群出现在对一维量子力学系统的描述中，尤其是在Stone–von Neumann 定理的语境下。更一般地，可以考虑与$n$维系统相关的海森堡群，最一般的情况则对应于任意的辛向量空间。
\subsection{三维情形}
在三维情形下，两个海森堡矩阵的乘积为：
$$
\begin{pmatrix}
1 & a & c \\
0 & 1 & b \\
0 & 0 & 1
\end{pmatrix}
\begin{pmatrix}
1 & a' & c' \\
0 & 1 & b' \\
0 & 0 & 1
\end{pmatrix}
=
\begin{pmatrix}
1 & a + a' & c + ab' + c' \\
0 & 1 & b + b' \\
0 & 0 & 1
\end{pmatrix}.~
$$
由其中的$ab'$项可以看出，该群是非阿贝尔群。

海森堡群的中性元是单位矩阵，其逆元由下式给出：
$$
\begin{pmatrix}
1 & a & c \\
0 & 1 & b \\
0 & 0 & 1
\end{pmatrix}^{-1}
=
\begin{pmatrix}
1 & -a & ab - c \\
0 & 1 & -b \\
0 & 0 & 1
\end{pmatrix}.~
$$
该群是二维仿射群$\mathrm{Aff}(2)$的一个子群：$\begin{pmatrix}1 & a & c \\0 & 1 & b \\0 & 0 & 1\end{pmatrix}$作用在$(\vec{x}, 1)$上时，对应的仿射变换为：$\begin{pmatrix}1 & a \\0 & 1\end{pmatrix} \vec{x} +\begin{pmatrix}c \\b\end{pmatrix}$.三维情形下有若干重要的例子。
\subsubsection{连续海森堡群}
若$a, b, c$为实数（在实数环$\mathbf{R}$中），则得到连续海森堡群$H_{3}(\mathbf{R})$。

它是一个3维幂零实李群。

除了作为实$3 \times 3$矩阵的表示外，连续海森堡群在函数空间中还有若干不同的表示。根据Stone–von Neumann 定理，在同构意义下，海森堡群$H$存在唯一的不变的不可约酉表示，其中其中心通过一个给定的非平凡特征作用。

这一表示有若干重要的实现方式（或模型）：在薛定谔模型中，海森堡群作用在平方可积函数空间上；在θ表示中，它作用在上半平面上的全纯函数空间上；其名称源自其与θ函数的联系。
\subsubsection{离散海森堡群}
\begin{figure}[ht]
\centering
\includegraphics[width=8cm]{./figures/ee8c83aac78d7ea1.png}
\caption{离散海森堡群的 Cayley 图的一部分，其生成元为文中所述的$x, y, z$（着色仅用于辅助视觉效果）。} \label{fig_HSBQ_1}
\end{figure}
如果$a, b, c$是整数（在环$\mathbf{Z}$ 中），则可以得到离散海森堡群$H_{3}(\mathbf{Z})$。这是一个非阿贝尔的幂零群。它有两个生成元：
$$
x = \begin{pmatrix}
1 & 1 & 0 \\
0 & 1 & 0 \\
0 & 0 & 1
\end{pmatrix}, \quad
y = \begin{pmatrix}
1 & 0 & 0 \\
0 & 1 & 1 \\
0 & 0 & 1
\end{pmatrix}~
$$
并且满足如下关系式：
$$
z = x y x^{-1} y^{-1}, \quad xz = zx, \quad yz = zy,~
$$
其中
$$
z = \begin{pmatrix}
1 & 0 & 1 \\
0 & 1 & 0 \\
0 & 0 & 1
\end{pmatrix}~
$$
是$H_{3}$的中心的生成元。（注意：矩阵$x, y, z$的逆元，只需把对角线上方的1替换成−1。）

根据Bass定理，它具有4阶的多项式增长率。

群中的任意元素都可以通过下式生成：
$$
\begin{pmatrix}
1 & a & c \\
0 & 1 & b \\
0 & 0 & 1
\end{pmatrix}
= y^{b} z^{c} x^{a}.~
$$
\subsubsection{模 $p$ 的海森堡群（当 $p$ 是奇素数时）}
如果取$a, b, c \in \mathbf{Z}/p\mathbf{Z}$，其中$p$为奇素数，那么就得到模 $p$ 的海森堡群。它是一个阶为$p^{3}$的群，生成元为$x, y$，并且满足以下关系：
$$
z = x y x^{-1} y^{-1}, \quad x^{p} = y^{p} = z^{p} = 1, \quad xz = zx, \quad yz = zy.~
$$
在奇素数阶有限域上的海森堡群的类似物称为额外特殊群，更准确地说，是指数为 $p$的额外特殊群。更一般地，如果一个群 $G$ 的导出子群包含在其中心 $Z$ 内，那么映射$G/Z \times G/Z \;\to\; Z$就是一个阿贝尔群上的斜对称双线性算子。

然而，若要求$G/Z$是有限向量空间，就必须要求$G$的 Frattini 子群包含在中心中；而若要求$Z$是$\mathbf{Z}/p\mathbf{Z}$上的一维向量空间，则必须要求 $Z$的阶为$p$。因此，如果 $G$ 不是阿贝尔群，那么$G$就是一个额外特殊群。如果$G$是额外特殊群，但其指数不是$p$，那么下面给出的关于辛向量空间$G/Z$的一般构造并不会得到与$G$同构的群。
\subsubsection{模 2 的海森堡群}
模2的海森堡群阶为8，并且同构于二面体群$D_{4}$（即正方形的对称群）。注意到如果
$$
x = \begin{pmatrix}
1 & 1 & 0 \\
0 & 1 & 0 \\
0 & 0 & 1
\end{pmatrix}, \quad
y = \begin{pmatrix}
1 & 0 & 0 \\
0 & 1 & 1 \\
0 & 0 & 1
\end{pmatrix},~
$$
那么
$$
xy = \begin{pmatrix}
1 & 1 & 1 \\
0 & 1 & 1 \\
0 & 0 & 1
\end{pmatrix}, \quad
yx = \begin{pmatrix}
1 & 1 & 0 \\
0 & 1 & 1 \\
0 & 0 & 1
\end{pmatrix}.~
$$
元素$x$ 和 $y$对应于两个相隔$45^\circ$的反射，而$xy$和$yx$对应于$90^\circ$的旋转。其余的反射是$xyx$与$yxy$，而$180^\circ$的旋转则由$xyxy \; (= yxyx)$给出。
\subsection{海森堡代数}
海森堡群$H$（在实数域上）的李代数$\mathfrak{h}$称为海森堡代数。[1]
它可以表示为如下形式的$3 \times 3$矩阵空间 [2]：
$$
\begin{pmatrix}
0 & a & c \\
0 & 0 & b \\
0 & 0 & 0
\end{pmatrix}, \quad a, b, c \in \mathbb{R}.~
$$
下列三个元素构成$\mathfrak{h}$的一组基：
$$
X = \begin{pmatrix}
0 & 1 & 0 \\
0 & 0 & 0 \\
0 & 0 & 0
\end{pmatrix}, \quad
Y = \begin{pmatrix}
0 & 0 & 0 \\
0 & 0 & 1 \\
0 & 0 & 0
\end{pmatrix}, \quad
Z = \begin{pmatrix}
0 & 0 & 1 \\
0 & 0 & 0 \\
0 & 0 & 0
\end{pmatrix}.~
$$
这些基元素满足以下对易关系：
$$
[X, Y] = Z, \quad [X, Z] = 0, \quad [Y, Z] = 0.~
$$
“海森堡群”这一名称的动机来自于上述关系，它们与量子力学中的正则对易关系形式相同：
$$
[\hat{x}, \hat{p}] = i \hbar I, \quad [\hat{x}, i \hbar I] = 0, \quad [\hat{p}, i \hbar I] = 0,~
$$
其中$\hat{x}$是位置算符，$\hat{p}$是动量算符，$\hbar$是普朗克常数。

海森堡群$H$具有一个特殊性质：指数映射是一个从李代数$\mathfrak{h}$到群$H$的一一对应的满射[3]：
$$
\exp \begin{pmatrix}
0 & a & c \\
0 & 0 & b \\
0 & 0 & 0
\end{pmatrix}
=
\begin{pmatrix}
1 & a & c + \tfrac{ab}{2} \\
0 & 1 & b \\
0 & 0 & 1
\end{pmatrix}.~
$$
\subsubsection{在共形场论中}
在共形场论中，“海森堡代数”一词被用来指代上述代数的无限维推广。它由元素$a_{n}, \quad n \in \mathbb{Z}$所张成，并满足以下对易关系：
$$
[a_{n}, a_{m}] = \delta_{n+m,0}.~
$$
在重新缩放之后，这实际上就是上述有限维海森堡代数的可数无穷多个拷贝。
\subsection{更高维度}
在更高维的欧几里得空间中，以及更一般的辛向量空间上，可以定义更一般的海森堡群$H_{2n+1}$。最简单的一般情形是维数为$2n+1$的实海森堡群，其中 $n \geq 1$ 为任意整数。作为矩阵群，$H_{2n+1}$（或者记为$H_{2n+1}(\mathbb{R})$，表示这是定义在实数域 $\mathbb{R}$上的海森堡群）被定义为一类$(n+2)\times(n+2)$的实矩阵群，其矩阵形式为：
$$
\begin{bmatrix}
1 & \mathbf{a} & c \\
\mathbf{0} & I_{n} & \mathbf{b} \\
0 & \mathbf{0} & 1
\end{bmatrix},~
$$
其中：\\
$\mathbf{a}$是长度为$n$的行向量；\\
$\mathbf{b}$是长度为$n$的列向量；\\
$I_{n}$ 是 $n \times n$的单位矩阵。
\subsubsection{群结构}
这确实是一个群，其乘法运算如下所示：
$$
\begin{bmatrix}
1 & \mathbf{a} & c \\
\mathbf{0} & I_{n} & \mathbf{b} \\
0 & \mathbf{0} & 1
\end{bmatrix}
\cdot
\begin{bmatrix}
1 & \mathbf{a}' & c' \\
\mathbf{0} & I_{n} & \mathbf{b}' \\
0 & \mathbf{0} & 1
\end{bmatrix}
=
\begin{bmatrix}
1 & \mathbf{a} + \mathbf{a}' & c + c' + \mathbf{a} \cdot \mathbf{b}' \\
\mathbf{0} & I_{n} & \mathbf{b} + \mathbf{b}' \\
0 & \mathbf{0} & 1
\end{bmatrix},~
$$
并且
$$
\begin{bmatrix}
1 & \mathbf{a} & c \\
\mathbf{0} & I_{n} & \mathbf{b} \\
0 & \mathbf{0} & 1
\end{bmatrix}
\cdot
\begin{bmatrix}
1 & -\mathbf{a} & -c + \mathbf{a} \cdot \mathbf{b} \\
\mathbf{0} & I_{n} & -\mathbf{b} \\
0 & \mathbf{0} & 1
\end{bmatrix}
=
\begin{bmatrix}
1 & \mathbf{0} & 0 \\
\mathbf{0} & I_{n} & \mathbf{0} \\
0 & \mathbf{0} & 1
\end{bmatrix}.~
$$
\subsubsection{李代数}
海森堡群是一个单连通李群，其李代数由如下形式的矩阵组成：
$$
\begin{bmatrix}
0 & \mathbf{a} & c \\
\mathbf{0} & 0_{n} & \mathbf{b} \\
0 & \mathbf{0} & 0
\end{bmatrix},~
$$
其中：\\
$\mathbf{a}$ 是长度为 $n$ 的行向量；\\
$\mathbf{b}$ 是长度为 $n$ 的列向量；\\
$0_{n}$ 是 $n \times n$ 的零矩阵。\\
设$e_{1}, \ldots, e_{n}$为$\mathbb{R}^{n}$的标准基，并定义：
$$
p_{i} =
\begin{bmatrix}
0 & e_{i}^{T} & 0 \\
\mathbf{0} & 0_{n} & \mathbf{0} \\
0 & \mathbf{0} & 0
\end{bmatrix}, \quad
q_{j} =
\begin{bmatrix}
0 & \mathbf{0} & 0 \\
\mathbf{0} & 0_{n} & e_{j} \\
0 & \mathbf{0} & 0
\end{bmatrix}, \quad
z =
\begin{bmatrix}
0 & \mathbf{0} & 1 \\
\mathbf{0} & 0_{n} & \mathbf{0} \\
0 & \mathbf{0} & 0
\end{bmatrix}.~
$$
由此得到的李代数满足以下正则对易关系：
$$
[p_{i}, q_{j}] = \delta_{ij} z, \quad [p_{i}, z] = 0, \quad [q_{j}, z] = 0,~
$$
其中$p_{1}, \ldots, p_{n}, q_{1}, \ldots, q_{n}, z$是代数的生成元。

特别地，$z$是海森堡李代数的中心元。需要注意的是，海森堡群的李代数是幂零的。
\subsubsection{指数映射}
设
$$
u = 
\begin{bmatrix}
0 & \mathbf{a} & c \\
\mathbf{0} & 0_{n} & \mathbf{b} \\
0 & \mathbf{0} & 0
\end{bmatrix},~
$$
它满足$u^{3} = 0_{n+2}$.指数映射计算为：
$$
\exp(u) = \sum_{k=0}^{\infty} \frac{1}{k!} u^{k} 
= I_{n+2} + u + \tfrac{1}{2} u^{2} 
= 
\begin{bmatrix}
1 & \mathbf{a} & c + \tfrac{1}{2}\mathbf{a} \cdot \mathbf{b} \\
\mathbf{0} & I_{n} & \mathbf{b} \\
0 & \mathbf{0} & 1
\end{bmatrix}.~
$$
任意幂零李代数的指数映射，都是该李代数与其唯一对应的连通、单连通李群之间的一个微分同胚。

除了涉及维度和李群的表述外，如果我们将$\mathbf{R}$替换为任意交换环$A$，上述讨论同样适用。相应的群记作$H_{n}(A)$。

在进一步假设素数 2 在环$A$中可逆的条件下，指数映射同样成立，因为它退化为有限项的求和，并且形式与上式相同（例如，$A$可以是奇素数$p$的剩余类环 $\mathbf{Z}/p\mathbf{Z}$，或任何特征为 0 的域）。
\subsection{表示论}
海森堡群的酉表示论相当简单——后来由Mackey理论推广——而且正是它在量子物理中被引入的动机，如下所述。

对于每个非零实数$\hbar$，可以定义海森堡群$H_{2n+1}$的一个不可约酉表示$\Pi_{\hbar}$，它作用在希尔伯特空间$L^{2}(\mathbb{R}^{n})$上，其定义公式为 [4]：
$$
\bigg[\Pi_{\hbar}
\begin{pmatrix}
1 & \mathbf{a} & c \\
0 & I_{n} & \mathbf{b} \\
0 & 0 & 1
\end{pmatrix}
\psi \bigg](\mathbf{x})
=
e^{i\hbar c} e^{i \mathbf{b}\cdot \mathbf{x}} \psi(\mathbf{x} + \hbar \mathbf{a})~
$$
这种表示称为薛定谔表示。它的动机来自于量子力学中位置算符与动量算符的指数形式的作用。参数$\mathbf{a}$表示在位置空间中的平移；参数$\mathbf{b}$表示在动量空间中的平移；参数$c$给出一个整体的相位因子。需要这个相位因子才能形成一个真正的算符群，因为位置空间的平移与动量空间的平移是不对易的。

关键结果是石–冯·诺伊曼定理，它表明：海森堡群的每一个（强连续的）不可约酉表示，只要其中心的作用是非平凡的，那么它必然与某个$\Pi_{\hbar}$（对应某个 $\hbar$）等价 [5]。换句话说，它们都等价于定义在$2n$维辛空间上的Weyl 代数（或 CCR 代数）。

由于海森堡群是$\mathbb{R}^{2n}$的一个一维中心扩张，所以它的不可约酉表示也可以被看作是$\mathbb{R}^{2n}$的不可约酉投射表示。从概念上讲，上面给出的表示就是经典相空间$\mathbb{R}^{2n}$上平移对称群的量子力学对应物。量子版本只是$\mathbb{R}^{2n}$的一个投射表示，这一点在经典层面就已经有所暗示。相空间中平移的哈密顿生成元是位置函数和动量函数。然而，它们在 Poisson 括号下的张成并不构成一个李代数，因为$\{x_{i}, p_{j}\} = \delta_{i,j}$.实际上，位置函数、动量函数再加上常数函数的张成，才在Poisson括号下构成一个李代数。这个李代数是交换李代数$\mathbb{R}^{2n}$的一个一维中心扩张，同构于海森堡群的李代数。
\subsubsection{在辛向量空间上的情形}
海森堡群的一般抽象形式可以由任意辛向量空间构造而成[6]。例如，设$(V, \omega)$是一个有限维实辛向量空间（因此$\omega$是$V$上的一个非退化斜对称双线性形式）。在$(V, \omega)$上的海森堡群$H(V)$（或简记为$V$）是集合 $V \times \mathbf{R}$，并赋予如下群运算：
$$
(v, t) \cdot (v', t') = \big(v + v', \; t + t' + \tfrac{1}{2}\omega(v, v')\big).~
$$
海森堡群是加法群$V$的一个中心扩张。因此存在如下精确序列：
$$
0 \;\longrightarrow\; \mathbf{R} \;\longrightarrow\; H(V) \;\longrightarrow\; V \;\longrightarrow\; 0.~
$$
任意辛向量空间都承认一个Darboux基$\{e_{j}, f_{k}\}_{1 \le j,k \le n}$，满足$\omega(e_{j}, f_{k}) = \delta_{jk}$，其中$V$的维数为 $2n$（必然是偶数）。在这个基的表示下，每个向量可以分解为：
$$
v = q^{a} e_{a} + p_{a} f^{a}.~
$$
这里的$q^{a}$与$p_{a}$是正则共轭坐标。

如果$\{e_{j}, f_{k}\}_{1 \leq j,k \leq n}$是$V$的一个Darboux基，那么取$\{E\}$作为$\mathbb{R}$ 的基，则$\{e_{j}, f_{k}, E\}_{1 \leq j,k \leq n}$就是$V \times \mathbb{R}$的相应基。此时，海森堡群中的一个向量可表示为：
$$
v = q^{a} e_{a} + p_{a} f^{a} + tE.~
$$
群运算变为：
$$
(p, q, t) \cdot (p', q', t') = \big(p + p',\; q + q',\; t + t' + \tfrac{1}{2}(pq' - p'q)\big).~
$$
由于海森堡群的底层流形是一个线性空间，因此李代数中的向量可以规范地与群中的向量对应。海森堡群的李代数由以下对易关系给出：
$$
[(v_{1}, t_{1}), (v_{2}, t_{2})] = \omega(v_{1}, v_{2}),~
$$
或者用 Darboux 基来写：
$$
[e_{a}, f^{b}] = \delta_{a}^{b},~
$$
并且所有其他对易子均为零。
